How much does retrieval actually help? Our experiments give a clear answer.
Table~\ref{tab:main_results} shows the full results across all ten
configurations, and Figure~\ref{fig:accuracy} makes the pattern immediately
visible: every RAG pipeline beats the no-RAG baseline.

To put those numbers in context: a random guesser on our three-class problem
would score 33.3\%. No-RAG sits barely above that (37.5\%), which confirms
that our benchmark is genuinely hard --- the model cannot rely on memorized
knowledge to answer correctly.

\subsection{Main Benchmark Results}

\begin{table*}[t]
\centering
\footnotesize
\caption{Results on the 40-case RAG-aware benchmark (N=40, three classes).
  Accuracy = exact-match label accuracy. Cost = total API cost for all 40
  evaluations. Best accuracy highlighted in bold.}
\label{tab:main_results}
\setlength{\tabcolsep}{8pt}
\begin{tabular}{llrrrr}
\toprule
Pipeline      & Model        & Accuracy        & Mean Score & Latency (s) & Cost (USD) \\
\midrule
BM25          & GPT-4o       & \textbf{47.5\%} & 70.6       & 2.36        & 0.144 \\
BM25          & GPT-4o-mini  & 45.0\%          & 80.2       & 2.15        & \textbf{0.008} \\
Dense         & GPT-4o       & 32.5\%          & 68.1       & 1.99        & 0.146 \\
Dense         & GPT-4o-mini  & 45.0\%          & 77.2       & 2.32        & \textbf{0.008} \\
Hybrid        & GPT-4o       & \textbf{47.5\%} & 69.9       & 1.82        & 0.145 \\
Hybrid        & GPT-4o-mini  & \textbf{47.5\%} & 77.6       & 1.71        & \textbf{0.008} \\
Multi-Query   & GPT-4o       & 45.0\%          & 71.4       & 1.94        & 0.145 \\
Multi-Query   & GPT-4o-mini  & 42.5\%          & 80.0       & 2.18        & \textbf{0.008} \\
No-RAG        & GPT-4o       & 37.5\%          & 74.1       & 4.06        & 0.060 \\
No-RAG        & GPT-4o-mini  & 37.5\%          & 87.8       & 1.90        & 0.003 \\
\midrule
Random chance & ---          & 33.3\%          & ---        & ---         & --- \\
\bottomrule
\end{tabular}
\end{table*}

\paragraph{Which pipeline wins?}
Hybrid and BM25 are the clear winners at \textbf{47.5\%} --- a
\textbf{+10 percentage point gain} over no-RAG.
Dense retrieval is the surprise loser: Dense + GPT-4o scores only 32.5\%,
actually \textit{below} the no-RAG baseline.
We think BM25's exact-match scoring handles specialized Turkish legal
terminology better than vector similarity in this domain.
Multi-Query performs in the middle, suggesting that expanding the query
sometimes helps but can also pull in off-topic passages.

\paragraph{Does the bigger model help?}
Mostly no --- and that is actually good news for deployment.
As Figure~\ref{fig:model_comparison} shows, GPT-4o-mini matches or beats
GPT-4o in most configurations.
For Hybrid retrieval, both tie at 47.5\% while GPT-4o costs
\textbf{19$\times$ more} per query (\$0.145 vs.\ \$0.008 for 40 cases).
Figure~\ref{fig:cost} shows the full cost-performance trade-off:
GPT-4o-mini points consistently appear in the high-accuracy, low-cost region.

\subsection{Where Models Struggle: Per-Class Breakdown}

The overall accuracy numbers tell only part of the story.
Table~\ref{tab:perclass} breaks down recall by compliance class, and
Figure~\ref{fig:heatmap} visualizes this as a heatmap.

\begin{table}[H]
\centering
\footnotesize
\caption{Per-class recall for selected configurations.
  Compliant: 13 cases. Partial: 14 cases. Non-compliant: 13 cases.}
\label{tab:perclass}
\begin{tabular}{llrrr}
\toprule
Pipeline    & Model        & Comp. & Partial & Non-Comp. \\
\midrule
Hybrid      & GPT-4o       & 46\%  & 93\%    & 0\% \\
Hybrid      & GPT-4o-mini  & 54\%  & 71\%    & 15\% \\
BM25        & GPT-4o       & 46\%  & 93\%    & 0\% \\
BM25        & GPT-4o-mini  & 62\%  & 57\%    & 15\% \\
No-RAG      & GPT-4o       & 31\%  & 71\%    & 8\% \\
No-RAG      & GPT-4o-mini  & 69\%  & 43\%    & 0\% \\
\bottomrule
\end{tabular}
\end{table}

The pattern in Figure~\ref{fig:heatmap} is striking.
Models are excellent at spotting \textit{partial} compliance (up to 93\%),
but almost completely miss \textit{non-compliant} documents ---
GPT-4o in the Hybrid configuration scores \textbf{0\%} on non-compliant cases.
Every time a document clearly violates a regulation, the model defaults
to ``partial'' instead.

This is not a quirk of our system --- it reflects a well-known tendency of
large language models to hedge rather than make strong negative
judgments \citep{Fazlija2024, Deshpande2024}.
For anyone deploying ComplianceAI in practice, this matters:
the system is good at flagging documents that are \textit{not clearly okay},
but human reviewers should make the final call on anything flagged as
non-compliant. Figure~\ref{fig:score_dist} further illustrates this
by showing how predicted compliance scores distribute across all pipelines.

\subsection{How Much Context Is Enough? Chunk Size Ablation}

One of the most practical decisions in building a RAG system is how to
chunk the source documents.
Table~\ref{tab:chunk_ablation} and Figure~\ref{fig:chunk_ablation} show
what happens as we sweep chunk sizes from 150 to 500 words.

\begin{table}[H]
\centering
\footnotesize
\caption{Chunk size ablation (BM25, GPT-4o-mini, Top-$K{=}5$).}
\label{tab:chunk_ablation}
\begin{tabular}{lrr}
\toprule
Chunk Size (words) & Accuracy & Latency (s) \\
\midrule
150                & 47.5\%   & 2.26 \\
250                & 50.0\%   & 1.95 \\
\textbf{350}       & \textbf{52.5\%} & 2.13 \\
500                & 40.0\%   & 2.11 \\
\bottomrule
\end{tabular}
\end{table}

The result is a clear inverted-U curve: accuracy climbs from 150 to 350
words, then drops sharply at 500.
\textbf{350 words is optimal at 52.5\%} --- 12 percentage points above
the 500-word setting.
YÖK regulation articles typically span 200--400 words, so a 350-word chunk
captures one complete article most of the time.
At 500 words, two articles get bundled together and the model receives
mixed signals.

\subsection{How Deep to Retrieve? Top-K Ablation}

Table~\ref{tab:topk_ablation} and Figure~\ref{fig:topk_ablation} show
the effect of retrieval depth $K$.

\begin{table}[H]
\centering
\footnotesize
\caption{Top-$K$ ablation (BM25, GPT-4o-mini, chunk size = 350).}
\label{tab:topk_ablation}
\begin{tabular}{lrr}
\toprule
Top-$K$ & Accuracy & Latency (s) \\
\midrule
1       & 50.0\%   & 1.85 \\
3       & 50.0\%   & 1.95 \\
\textbf{5}  & \textbf{52.5\%} & 2.04 \\
10      & 50.0\%   & 2.07 \\
\bottomrule
\end{tabular}
\end{table}

$K{=}5$ is optimal at 52.5\%, while other values plateau at 50.0\%.
Five chunks are enough to cover most compliance questions, which typically
involve one or two regulation articles plus some surrounding context.
Beyond five, extra chunks add latency but not accuracy.

\subsection{Response Time in Practice}

Figure~\ref{fig:latency} summarizes mean latency across all configurations.
The fastest RAG option is \textbf{Hybrid + GPT-4o-mini at 1.71 seconds}
per request --- fast enough for interactive use.
No-RAG + GPT-4o is counterintuitively the \textit{slowest} (4.06 seconds),
because without retrieval the full document is included in the prompt,
making inputs much longer.

\subsection{Key Takeaways}

Our experiments lead to three practical recommendations:

\begin{enumerate}
  \item \textbf{Retrieve. Always.}
    Even simple BM25 adds 10 percentage points over no-RAG.
    Hybrid is the most reliable choice.

  \item \textbf{Use GPT-4o-mini.}
    It matches GPT-4o accuracy at 19$\times$ lower cost.
    For a compliance tool processing hundreds of documents per year,
    this is a deciding factor.

  \item \textbf{Keep humans in the loop for non-compliant verdicts.}
    Models systematically miss violations. Use ComplianceAI as a
    first-pass filter, not a final judge.
\end{enumerate}
